\chapter{Tab 1 · Load files: reading and decoding}\label{ch:load}

\figui{ui_start.png}{The page before any file is read.}{ui_start}

\section{File intake}
Files arrive by drag and drop or the file dialog. \texttt{intake()} sorts them by extension (a \texttt{.zip} is opened and searched
for \texttt{.ixo}, \texttt{.eds} and \texttt{.edt} entries), decodes .ixo instrument XML text once, and lists what will be read.
\texttt{readStaged()} then decodes each file, runs the control-panel extraction while the bytes are still in memory and
opens the Control panel.

\par\medskip\noindent{\small\textbf{Listing}~\texttt{intake()}\hfill\textcolor{gray}{HTML lines 9993--10062}}\par\nopagebreak
\lstinputlisting[style=vstyle,language=JavaScript,firstnumber=9993]{code/intake.js}


\par\medskip\noindent{\small\textbf{Listing}~\texttt{readStaged()}\hfill\textcolor{gray}{HTML lines 10063--10073}}\par\nopagebreak
\lstinputlisting[style=vstyle,language=JavaScript,firstnumber=10063]{code/readStaged.js}



\figui[0.9]{ui_load.png}{Load tab after reading: the synthetic data set notice and the \emph{What was read} table (one row per run).}{ui_load}

\section{.ixo instrument (\texttt{.ixo})}
An \texttt{.ixo} file is XML. Objects (\texttt{<obj>}) carry properties (\texttt{<prop name="\dots">}); binary members are base64.
The decoder reads the plate model, subsets, protocol, analyses, results and the raw acquisitions.

\par\medskip\noindent{\small\textbf{Listing}~\texttt{decodeIxo() — opening part}\hfill\textcolor{gray}{HTML lines 2687--3029}}\par\nopagebreak
\lstinputlisting[style=vstyle,language=JavaScript,firstnumber=2687]{code/decodeIxo.js}



\subsection*{Container integrity}
The file stores a checksum of its content. \texttt{ixoIntegrity()} recomputes it; a mismatch means the XML was edited
outside the instrument software and is reported as an error in the forensic log.

\par\medskip\noindent{\small\textbf{Listing}~\texttt{ixoIntegrity()}\hfill\textcolor{gray}{HTML lines 3742--3756}}\par\nopagebreak
\lstinputlisting[style=vstyle,language=JavaScript,firstnumber=3742]{code/ixoIntegrity.js}



\subsection*{Block temperature log (\texttt{TemperatureLog})}
The block temperature is logged about every 0.41\,s for the whole run. Two arrays are stored: \texttt{Times} with the magic
\texttt{LARZ} (uint32 milliseconds) and \texttt{Temps} with \texttt{DARZ} (float64 °C). After a short header the values are
either raw or compressed with LZSS (4\,096-byte ring buffer, matches of 3--18 bytes, one flag bit per item).

\begin{infobox}
The header length varies between files (12 or 13 bytes were found). Reading from the wrong offset produced temperatures such as
$10^{102}$\,°C for one interrupted run, which distorted every thermal metric of that run. \texttt{ccDecodeArz()} therefore
tries every plausible start and keeps the one whose \emph{whole} array is physically plausible (temperatures $-50$\dots$200$\,°C,
times never decreasing); \texttt{ccIxoTemperatureLog()} drops any point that is still implausible.
\end{infobox}

\par\medskip\noindent{\small\textbf{Listing}~\texttt{ccLzss()}\hfill\textcolor{gray}{HTML lines 5545--5559}}\par\nopagebreak
\lstinputlisting[style=vstyle,language=JavaScript,firstnumber=5545]{code/ccLzss.js}


\par\medskip\noindent{\small\textbf{Listing}~\texttt{ccDecodeArz()}\hfill\textcolor{gray}{HTML lines 5560--5573}}\par\nopagebreak
\lstinputlisting[style=vstyle,language=JavaScript,firstnumber=5560]{code/ccDecodeArz.js}


\par\medskip\noindent{\small\textbf{Listing}~\texttt{ccIxoTemperatureLog()}\hfill\textcolor{gray}{HTML lines 5574--5582}}\par\nopagebreak
\lstinputlisting[style=vstyle,language=JavaScript,firstnumber=5574]{code/ccIxoTemperatureLog.js}



\subsection*{Fluorescence acquisitions (\texttt{AcquisitionStore})}
Every fluorescence reading is a record with its time, the block temperature at that moment (hexadecimal big-endian double),
the channel, the automatic integration time, a validity flag and the lamp reference value. The store is deflate-compressed XML.

\par\medskip\noindent{\small\textbf{Listing}~\texttt{ccIxoAcquisitions()}\hfill\textcolor{gray}{HTML lines 5585--5592}}\par\nopagebreak
\lstinputlisting[style=vstyle,language=JavaScript,firstnumber=5585]{code/ccIxoAcquisitions.js}



\section{.eds/.edt instrument (\texttt{.eds}, \texttt{.edt})}
An \texttt{.eds} file is a ZIP container. \texttt{parseEds()} checks the CRC of every entry and the tamper value written by the
instrument software, then reads the experiment description, plate, protocol, analysis results, multicomponent and raw data,
calibrations, the per-image \texttt{quant} files and the instrument's \texttt{messages.log}. \texttt{decodeEds()} maps the result
onto the common run model.

\par\medskip\noindent{\small\textbf{Listing}~\texttt{parseEds() — opening part}\hfill\textcolor{gray}{HTML lines 3181--3387}}\par\nopagebreak
\lstinputlisting[style=vstyle,language=JavaScript,firstnumber=3181]{code/parseEds.js}


\par\medskip\noindent{\small\textbf{Listing}~\texttt{decodeEds() — opening part}\hfill\textcolor{gray}{HTML lines 3611--3739}}\par\nopagebreak
\lstinputlisting[style=vstyle,language=JavaScript,firstnumber=3611]{code/decodeEds.js}



\subsection*{Instrument log and Ct re-derivation}
\texttt{parseLog()} reads the once-per-second \texttt{Temperature} and \texttt{LEDStatus} lines. \texttt{recalcCt()} re-derives
every Ct from the stored $\Delta$Rn and threshold, so a stored Ct that does not match its own curve can be flagged.

\par\medskip\noindent{\small\textbf{Listing}~\texttt{parseLog()}\hfill\textcolor{gray}{HTML lines 3397--3413}}\par\nopagebreak
\lstinputlisting[style=vstyle,language=JavaScript,firstnumber=3397]{code/parseLog.js}


\par\medskip\noindent{\small\textbf{Listing}~\texttt{recalcCt()}\hfill\textcolor{gray}{HTML lines 3391--3395}}\par\nopagebreak
\lstinputlisting[style=vstyle,language=JavaScript,firstnumber=3391]{code/recalcCt.js}



\texttt{ccQsLog()} makes one more pass over the log for the timing and mechanics used by the control panel: command reply times,
camera and filter-wheel moves, encoder offsets, heated-cover drive, scheduler lateness, image processing, predicted and actual
run time, errors and the lifetime counters of the block.

\par\medskip\noindent{\small\textbf{Listing}~\texttt{ccQsLog()}\hfill\textcolor{gray}{HTML lines 5712--5752}}\par\nopagebreak
\lstinputlisting[style=vstyle,language=JavaScript,firstnumber=5712]{code/ccQsLog.js}



\section{What each file type carries}
\begin{center}\small
\begin{tabular}{@{}p{4.6cm}p{4.9cm}p{5.6cm}@{}}\toprule
Information & .ixo instrument (\texttt{.ixo}) & .eds/.edt instrument (\texttt{.eds})\\\midrule
Block temperature & \texttt{TemperatureLog}, 0.41\,s, block & \texttt{messages.log}, 1\,s, modelled sample temperature, 6 zones, cover, heat sink\\
Temperature at each reading & every acquisition record & every image (6 zones)\\
Light source & lamp reference channel per reading & LED temperature, current, voltage, junction\\
Exposure & automatic integration time per reading & requested exposure per image\\
Mechanics & channel switch timing & filter wheel, encoder, heated cover\\
Embedded computer & -- & command replies, 1-s timer, scheduler, image processing\\
Calibrations & -- & uniformity, background, ROI files\\
Lifetime counters & -- & block cycles, degrees climbed\\
Not in the file & lamp hours, plate loading, error log (Instrument Problem Report) & drawer and plate insertion, LED hours\\\bottomrule
\end{tabular}
\end{center}

\section{Extraction for the control panel}
\texttt{ccExtractOnLoad()} runs once per file during reading and stores the extracted instrument numbers on the run, together
with a thinned copy of the temperature log for display.

\par\medskip\noindent{\small\textbf{Listing}~\texttt{ccExtractOnLoad()}\hfill\textcolor{gray}{HTML lines 5901--5920}}\par\nopagebreak
\lstinputlisting[style=vstyle,language=JavaScript,firstnumber=5901]{code/ccExtractOnLoad.js}




